\chapter{刑法论}

% Source PDF physical page 13; see audit/source_coverage.json.
\begin{EditionOpeningPanel}
\InsertDiagram{figures/ch01_overview.tex}
\end{EditionOpeningPanel}

\section{刑法的解释}

[问题]狗蛋给室友的饭盆撒尿，是否属于故意毁坏财物罪中的“毁坏”？\footnote{“毁坏”包括物理性毁坏和价值性毁坏。室友没法用饭盆了，饭盆的价值丧失了。所以，狗蛋的行为属于毁坏财物。}

\InsertDiagram{figures/ch01_research.tex}

% Source PDF physical page 14; see audit/source_coverage.json.
刑法解释学又被称为刑法教义学（释义学），意思是将法条规定当作教义一样去解释，培养人们对法律的信仰。“法律不是被嘲笑的对象”。这也是法考命题的基本理念。基于这一理念，法考考查的能力是，如何解释法律、适用法律，妥当处理案件，而不是如何批判法律、修改法律。因此，刑法的解释方法成为历年考试重点。

\CJKunderwave{第一，解释对象}。解释的对象是现行刑法规范。这包括两个：一是刑法典，1997年10月1日生效；刑法修正案是对刑法典作的小修小补，至今已有12个刑法修正案。二是一部单行刑法。\footnote{单行刑法，是指国家在刑法典之外单独针对某一类犯罪规定的刑法规范。目前典型的单行刑法只有\CJKunderwave{一个}：1998年12月29日全国人大常委会颁布的《关于惩治骗购外汇、逃汇和非法买卖外汇犯罪的决定》。提示：附属刑法，是指附带规定在经济法、行政法等非刑事法律中的罪刑规范。有些国家的法律有此做法，例如，德国、日本将许多交通领域的犯罪规定在道路交通法中。在我国，只有刑法典和单行刑法有权规定犯罪和刑罚，因此我国目前没有附属刑法。}

\CJKunderwave{第二，解释效力}。全国人大常委会作的立法解释\ensuremath{>}最高法院、最高检察院作的司法解释。学者作的学理解释没有法律效力。

\CJKunderwave{第三，解释方法}。这包括解释技巧和解释理由。\CJKunderwave{解释技巧负责生产结论，类似于生产}\CJKunderwave{部门；解释理由负责为结论提供论证理由，类似于质检部门}。

\subsection*{一、解释技巧}

对一个条文术语的含义，需要通过解释技巧得出解释结论。常考的解释技巧有：

\subsubsection*{（一）平义解释}

这是指按照该用语最平常的字面含义来解释，既不扩大含义范围，也不缩小含义范围，即所谓的看山是山，看水是水。例如（2009年第1题），将盗窃罪对象的“公私财物”解释为“他人的财物”，属于平义解释。“公私财物”意指不是自己的财物，而是“他人的财物”，这里的“他人”既包括自然人（私人），也包括单位（公家）；反正只要不是自己的财物，就能成为盗窃罪的对象。这种解释既没有扩大，也没有缩小盗窃罪的对象范围，因此属于平义解释。

\subsubsection*{（二）扩大解释（又称扩张解释）}

这是指对用语解释后的含义大于字面含义，但该含义仍处在该用语可能的含义范围内。例如，将侵犯通信自由罪中的“信件”解释为包含电子邮件，属于扩大解释。

\subsubsection*{（三）缩小解释（又称限制解释）}

这是指对用语解释后的含义小于字面含义，是字面含义的一个子集。例如（2018年试题），“为境外窃取、刺探国家情报罪”是危害国家安全的犯罪，因此，其中的“情报”不是泛指一切情报（如商业情报、高科技情报等），而应缩小解释为关系到国家安全的情报。

\subsubsection*{（四）类推解释}

这是指将不符合法律规定的情形解释为符合法律规定的情形。例如，将强奸罪中的“妇女”解释为包括男子，就属于类推解释。

1.性质。刑法禁止类推解释。无论立法解释，还是司法解释，均禁止类推解释。

2.例外。禁止类推解释是原则，但例外的，不禁止有利于被告人的类推解释。这是因为，禁止类推解释是为了保障人权，而有利于被告人的类推解释也是为了保障人权。

% Source PDF physical page 15; see audit/source_coverage.json.
3.区分。\CJKunderwave{类推解释与扩大解释的区分标准}（五星级考点）：第一，扩大解释得出的结论，在词义射程范围内；类推解释得出的结论，在词义射程范围外。第二，扩大解释得出的结论，没有超出国民的预测可能性；类推解释得出的结论，明显超出国民的预测可能性。简单讲：扩大解释，就是扯得还行；类推解释，就是扯得太远。

\InsertDiagram{figures/manual_p015_1.tex}

[总结]易考情形：\footnote{参见张明楷：《罪刑法定与刑法解释》，北京大学出版社2010年版。张明楷：《刑法分则的解释原理》（上、下册）（第二版），中国人民大学出版社2011年版。}

（1）将遗弃罪中的“负有扶养义务”的人解释为既包括家庭成员，也包括负有扶养义务的其他人，属于扩大解释。

（2）将侵犯通信自由罪中的“信件”解释为包含电子邮件，属于扩大解释。

（3）将“金融机构”解释为包含使用中的运钞车、自动取款机，属于扩大解释。（2025年试题）

（4）将走私弹药罪中的“弹药”解释为包含弹壳，属于扩大解释。这是司法解释的规定。

（5）将组织卖淫罪中的“组织他人卖淫”解释为包含男性向女性卖淫，属于扩大解释。

[巩固练习]题1，将信用卡诈骗罪中的“信用卡”解释为包含借记卡（储蓄卡），属于扩大解释。这是立法解释的规定（2009年第1题）。

题2，将破坏交通工具罪中的“汽车”解释为包含大型拖拉机，属于扩大解释（2013年第3题）；将劫持汽车罪中的“汽车”解释为包含火车、地铁，属于类推解释。如果劫持火车、地铁，可定破坏交通工具罪。

题3，将盗窃骨灰的行为认定为盗窃“尸体”，属于类推解释（2014年第3题）。

题4，将非法制造大炮解释为非法制造枪支罪，属于扩大解释（2018年试题）。

题5，将传播性病罪中的“性病”解释为包括艾滋病，属于扩大解释。这是司法解释的规定。

题6，将传播淫秽物品罪中的“淫秽物品”解释为包括“在网上直播淫秽表演”，属于扩大解释。

题7，将故意毁坏财物罪中的“毁坏”解释为，不仅包括物理性毁坏，还包括功能性毁坏，属于扩大解释。

\subsubsection*{（五）反义解释}

反义解释也称反对解释，是指根据用语的正面表述，推导出其反面含义，也即从“A”推导出“-A”。例如，16周岁及以上的人要对盗窃罪负刑事责任。据此，不满16周岁的人便不需对盗窃罪负刑事责任。（2020年试题）

% Source PDF physical page 16; see audit/source_coverage.json.
\subsection*{二、解释理由}

在通过解释技巧对一个条文术语得出一个结论后，必须提供理由，论证解释的合理性。常考的解释理由有：

\subsubsection*{（一）文理解释}

这是指根据文法、语法等来论证解释后的含义是否属于刑法用语可能具有的含义。简而言之，需要考查解释后的含义在文理上是否讲得通。

例1（2015年第51题），下面说法是否正确：刑法规定“以暴力、胁迫或者其他手段强奸妇女的”构成强奸罪。按照文理解释，可将丈夫强行与妻子性交的行为解释为“强奸妇女”。答案：说法正确。从文理上看，刑法第236条并没有将婚内强奸排除在强奸罪之外。所谓婚内强奸是否构成强奸罪的问题，主要是证据的证明问题，因为是夫妻关系，在证据上比较难证明强奸事实。一旦有证据证明，丈夫强奸了妻子，按照刑法第236条规定，应以强奸罪论处。

例2（2019年试题），招摇撞骗罪是指冒充国家机关工作人员招摇撞骗。将副乡长冒充市长招摇撞骗，解释为“冒充”国家机关工作人员招摇撞骗，符合文理解释。

\subsubsection*{（二）体系解释}

这是指根据体系逻辑来论证解释后的含义在刑法体系中是否协调合理。例如（2006年第20题），从文义上看，“伪造”可以包含“变造”，但是我国刑法在伪造货币罪之外又规定了变造货币罪，那么伪造货币罪中的“伪造”就不能包含“变造”。

1.“\CJKunderwave{同一用语的含义相对化}”（一词多义）。体系解释并不意味着同一用语在不同条文中需要保持同一含义。相反，同一用语在不同语境（不同条文、不同部门法）中可以保持不同含义。

例1（2015年第51题），“胁迫”既是抢劫罪的手段，也是强奸罪的手段，但在这两罪中的含义范围有所不同。抢劫罪的“胁迫”要求以暴力相胁迫，但强奸罪的“胁迫”可以是以非暴力相胁迫。例如，甲对乙声称：“不给钱，就曝光你的裸照！”甲不构成抢劫罪，而构成敲诈勒索罪；而甲对乙声称：“不答应跟我睡觉，就曝光你的裸照！”甲构成强奸罪。因此，强奸罪中“胁迫”的含义范围大于抢劫罪中“胁迫”的含义范围。\CJKunderwave{强奸罪}\CJKunderwave{中的“胁迫”=抢劫罪中的“胁迫”+敲诈勒索罪中的“胁迫”}。

例2（2013年第3题），“刑法分则有不少条文并列规定了‘伪造’与‘变造’，但不排除在其他一些条文中将‘变造’解释为‘伪造’的一种表现形式”。该说法是否正确？答案：正确。例如，伪造货币罪和变造货币罪，二者是并列关系，不是包含关系，因此，这里的“伪造”不能包含“变造”。又如，刑法规定了“帮助伪造证据罪”，没有规定“帮助变造证据罪”，因此“帮助伪造证据罪”中的“伪造”可以包括“变造”。可见，“伪造”在有些罪名中不能包括“变造”，在有些罪名中可以包括“变造”。

[巩固练习］题1（2016年第51题），强制猥亵、侮辱罪与侮辱罪，二者中的“侮辱”含义并不相同。前者是指侵犯人的性羞耻心。例如，强行脱妇女衣服。后者是指侵犯人的名誉。例如，往妇女身上泼洒粪便。

题2（2019年试题），传播淫秽物品牟利罪与传播性病罪，二者中的“传播”含义并不相同。前者是指使不特定人知晓内容。后者是指传染。

% Source PDF physical page 17; see audit/source_coverage.json.
题3，甲开设洗浴中心，只提供色情按摩服务，不提供性交服务。在行政法（治安管理处罚法）上，甲属于组织卖淫活动。但在刑法上甲不构成组织卖淫罪。刑法上的组织“卖淫”，含义范围要窄一些，要求提供性交服务。

2.“不同用语的含义同一化”（多词一义）。刑法中几个不同的用语也可以保持同一个含义。

例1，刑法条文中的“出售”“销售”“倒卖”“贩卖”的含义相同，是指有偿转让所有权。

例2，刑法条文中的“恐吓”“胁迫”“威胁”“敲诈”的含义相同，是指以恶害相通告，使对方产生恐惧心理。

3.同类解释规则。这是指立法者在描述罪状时经常使用“例示法”，也即先列举几个例子，然后用“等”“及其他”来概括兜底。对这些兜底规定的含义范围不能随意扩大，而应先总结所列举的例子的共同特征，然后用该共同特征来解释兜底规定的含义（2014年第3题）。

例如，抢劫罪规定的行为方式是“暴力、胁迫或者其他方法”。这里的“暴力、胁迫”是指压制被害人反抗的方法。因此，这里的“其他方法”是指其他的压制被害人反抗的方法，如将被害人灌醉弄昏，而非泛指其他一切方法。

\subsubsection*{（三）当然解释}

这是指在论证无罪时“举重以明轻”（重的行为都无罪，轻的行为当然更应无罪），在论证有罪时“举轻以明重”（轻的行为都是犯罪，重的行为当然更应是犯罪）。例1，如果抢劫不构成犯罪，那么盗窃更不构成犯罪。例2，如果盗窃构成犯罪，那么抢劫更构成犯罪。

[注意］当然解释所比较的两个事项必须是性质相同、程度不同的关系。如果性质不同，不能进行当然解释的推理。

例1，可以主张“强制猥亵都是犯罪，强奸更应是犯罪”，但不能主张“强制猥亵都是犯罪，抢劫更应是犯罪”。

例2（2015年第51题），“既然将为了自已饲养而抢劫他人宠物的行为认定为抢劫罪，那么，根据当然解释，对为了自已收养而抢劫他人婴儿的行为更应认定为抢劫罪，否则会导致罪刑不均衡”。该说法是否正确？答案：错误。有观点可能会如此比较，也符合朴素的正义感。但是，从刑法解释学角度看，宠物与婴儿是性质不同的两个概念，体现不同的法益，不能进行轻重比较，不能进行当然解释。例如，有观点将非法买卖国家保护动物（大熊猫、金丝猴）与买卖妇女、儿童进行比较，认为对后者应该处罚更重一些。但是，这两种概念属于性质不同的概念，无法进行当然解释。（2025年试题）

例3，《刑法》第329条规定了抢夺、窃取国有档案罪，但没有规定抢劫国有档案罪。然而，抢夺都是犯罪，抢劫更应是犯罪。抢劫行为可以评价为抢夺行为。所以，抢劫国有档案的行为可以按照抢夺国有档案罪论处。（2024年试题）

\subsubsection*{（四）目的解释}

这是指根据刑法的保护目的为解释的结论提供理由。

例如，乙为官清廉但不被提拔，心生一计，教唆甲诬告自已贪污，甲便诬告乙犯贪污罪。司法机关经过一番调查，发现乙没有贪污事实。甲是否构成诬告陷害罪？对此需要根据

% Source PDF physical page 18; see audit/source_coverage.json.
该罪的保护目的加以认定。（1）少数说认为，该罪的保护目的是司法机关的正常司法活动。甲破坏了司法机关的正常司法活动，浪费了司法资源，构成该罪。（2）多数说认为，该罪被规定在人身犯罪一章，保护目的是人身权利。乙同意甲诬告自己，表明乙放弃了对自己人身权利的保护，那么刑法便不需要保护乙的人身权利。所以，甲不构成诬告陷害罪。

[注意]\CJKunderwave{观点展示问题}

\CJKunderwave{本书主张客观题与主观题一体化复习}。刑法主观题经常考查观点展示。这是因为，刑法学不是自然科学，许多情况下没有唯一正确答案，存在多种观点。法考的任务是筛选和培养实务人士。实务人士比的是谁的理由更充分，更有说服力。实务人士的看家技能是说理论证。

1.对于不同观点，形成多数说与少数说的，本书一一标明哪个是多数说，哪个是少数说。\CJKunderwave{标出来的用途在于，如果问唯一正确答案，则按照多数说作答}。

2.有些不同观点，没有形成多数说与少数说，是平分秋色，对此不会考“唯一正确答案是什么”，但会考“\CJKunderwave{前提与结论之间的逻辑推导是否自洽}”。

3.观点展示是指命题组的观点展示，不是某个命题组成员个人的观点展示。当然，学术地位影响话语权权重。

4.法考真题素材更多来自实务案例，但是答案不一定与判决结论保持一致，因为有些判决结论是不妥当的，是老百姓不满意的，法考要发挥纠偏批判功能。

5.观点展示有两种，一是学术型观点展示，这个数量较少，可以背下来。二是实务型观点展示，也即法庭上控辩双方总是持不同观点，这种观点展示数量众多，没法背下来，也不需要背下来，主要是掌握一种逻辑思维和论证方法。法考对两种观点展示都考查，对实务型观点展示考查得很多。

\subsection*{三、解释技巧与解释理由的关系}

\subsubsection*{（一） 二者区别}

1.解释技巧是生产结论（生产部门）。生产结论的生产线只能选择一条，如果对一个词语进行扩大解释，就不能同时进行缩小解释（限制解释）。解释技巧之间的关系是排斥关系。

2.解释理由是论证结论的合理性（质检部门）。为结论提供论证理由，理由越多越好。解释理由之间的关系是并存关系。

例1，将组织卖淫罪中的“组织他人卖淫”解释为包括男性向女性卖淫。这是生产结论，属于扩大解释，此时就不可能是缩小解释。为该结论提供的理由可以多种多样，例如文理解释、体系解释、目的解释等。

例2（2016年第51题），“对甲法条中的‘暴力’作扩大解释时，就不可能同时再作限制解释，但这并不意味着对乙法条中的‘暴力’也须作扩大解释”。该说法是否正确？答案：正确。第一，扩大解释与限制解释（缩小解释）是排斥关系。第二，从体系角度看，一个词语在不同条文可以保持不同含义，也即一词多义。

% Source PDF physical page 19; see audit/source_coverage.json.
\subsubsection*{（二）二者关系}

解释技巧和解释理由之间是相辅相成、相互制约关系。前者负责生产结论，后者负责论证结论的合理性。通过解释技巧生产出结论，但论证理由不充分，例如，不符合目的解释或体系解释等，则不得采纳。反过来，根据解释理由论证某个结论很合理，但是生产该结论的手段是类推解释，则不得采纳。

例1（2016年第51题），“对刑法分则条文的解释，必须同时符合两个要求：一是不能超出刑法用语可能具有的含义，二是必须符合分则条文的目的。”该说法是否正确？答案：正确。第一句话是指，在解释技巧上，不能类推解释。第二句话是指，在解释理由上，要符合目的解释。二者相辅相成。

例2（2020年试题），“一个解释结论符合当然解释，但不符合刑法文字的可能含义，该结论符合罪刑法定原则。”该说法是否正确？答案：错误。一个结论虽然符合当然解释，但是在解释技巧上属于类推解释，不得采纳，否则违反罪刑法定原则。例如，刑法将“冒充军警人员抢劫”规定为抢劫罪的法定刑升格条件，加重处罚。根据当然解释，真警察实施抢劫，更应被加重处罚。但是，该结论难以从“冒充军警人员抢劫”中生产出来。若将“冒充军警人员抢劫”解释为包括真警察抢劫，应属于类推解释（多数说），故这种解释结论无法被采纳。\footnote{提示：文理解释与论理解释的关系。法理学的传统理论认为，解释方法分为“文理解释”与“论理解释”，将二者并列，认为二者是对立排斥关系。刑法新理论认为，“文理解释”和“论理解释”不是对立排斥关系，而是交叉重合关系。“论理解释”是一个笼统概念，既包含解释理由（如当然解释、体系解释等），也包含解释技巧（如扩大解释、缩小解释）；而文理解释是一种解释理由。由于解释理由与解释技巧是相辅相成关系，因此文理解释与论理解释不是对立关系。对一个刑法用语，既可以作文理解释，也可以作论理解释。这种新理论的立场是法考刑法的官方立场。有真题为证（2011年第51题）。

此外，学术界关于刑法的解释，还有客观解释与主观解释之争、形式解释与实质解释之争。由于法考不考这些概念，所以本书不予介绍。}

\OriginalPanel{典型真题\footnotemark}\footnotetext{本书选取的真题是代表命题方向的真题，主要用于考生在学习完精讲内容后及时进行检验，所以在理解答案时可以结合前文的精讲内容。更详细的答案解析可以参见与本书配套的《真金题卷》一书。}

关于刑法的解释，下列说法正确的是？（2019年回忆版金题）\footnote{【答案】A项，传播淫秽物品罪的“传播”是指让不特定人知晓淫秽物品的内容。传播性病罪中的“传播”是指通过性器官接触传染性病。二者含义不同。A项说法错误。B项，根据司法解释规定，B项说法正确。C项，冒充国家机关工作人员的情形包括：一是非国家机关工作人员冒充国家机关工作人员；二是此种国家机关工作人员冒充彼种国家机关工作人员，如行政机关工作人员冒充司法机关工作人员；三是职务低的国家机关工作人员冒充职务高的国家机关工作人员，当然也包括相反情形。C项说法错误。D项，虐待被监管人罪中的“体罚虐待”，是指体罚或者虐待，而非仅限于体罚式虐待。D项说法正确。本题答案：BD。}

A.按照体系解释，传播淫秽物品罪与传播性病罪的“传播”含义一致

B.依据论理解释，倒卖文物罪中的“倒卖”是指以牟利为目的，出售或为出售而购买国家禁止经营的文物

C.招摇撞骗罪是指冒充国家机关工作人员招摇撞骗。将副乡长冒充市长招摇撞骗解释为“冒充”国家机关工作人员招摇撞骗，不符合文理解释

D.虐待被监管人罪是指殴打或体罚虐待被监管人。将其中的“体罚虐待”解释为“体罚或者虐待”，符合文理解释

% Source PDF physical page 20; see audit/source_coverage.json.
\section{刑法的原则}

[问题]在拍卖会上，狗蛋组织多人串通拍卖。我国刑法没有规定串通拍卖罪，但规定有串通投标罪。能否给狗蛋定串通投标罪？（串通拍卖案，2020年试题）

刑法的原则包括三项：罪刑法定原则，罪刑相适应原则，刑法面前人人平等原则。法考很少考查第三项原则，故省略。

\subsection*{一、罪刑法定原则}

第3条\footnote{如果无特别提示，本书所列“第x条”均是指刑法的条文，均用阿拉伯数字表示。对于司法解释的条文，本书会特别提示。}法律明文规定为犯罪行为的，依照法律定罪处刑；法律没有明文规定为犯罪行为的，不得定罪处刑。

罪刑法定原则的含义是：法无明文规定不为罪，法无明文规定不处罚。这是刑法的铁则。

[注意］上述第3条第一句话并不是关于罪刑法定原则的规定，其旨意是防止司法人员将有罪行为作无罪处理，是在强调法益保护机能。第二句话才是关于罪刑法定原则的规定，是在强调人权保障机能。在保护法益与保障人权这个天平上，罪刑法定原则是用来保障人权的，而非保护法益。

\subsubsection*{（一）思想基础}

1.\CJKunderwave{民主主义}，意指只有民主产生的立法机关才有权规定犯罪与刑罚。

2.\CJKunderwave{自由主义}，也称为尊重人权主义，意指法律应当具有预测可能性，使国民知晓什么是犯罪，什么不是犯罪，以保障个人的行动自由。这种思想也称为国民预测可能性原理。

\subsubsection*{（二）基本内容}

1.\CJKunderwave{成文}的罪刑法定

（1）法律主义。这是指只有立法机关制定的法律才有权规定犯罪和刑罚，行政机关制定的法规等无权规定犯罪和刑罚，即罪刑法定中的“法”不包括行政法规等，也即国务院无权制定刑法。最高人民法院、最高人民检察院也无权制定刑法，只能解释刑法。法律主义是民主主义的体现。（2006年第1题）

（2）禁止习惯法。虽然习惯法体现民意，但因为不成文，缺乏明确性，违反了预测可能性原理，所以应当被禁止。（2006年第1题）

2.\CJKunderwave{事前}的罪刑法定

这是指禁止溯及既往（或禁止事后法）。刑法如果溯及既往，便违反了预测可能性原理。例如，我今天于件事，按照今天的法律，我无罪，但明天出台新的法律，规定这种行为构成犯罪，然后依照这个事后的法律给我定罪，这肯定不合理。

[注意］刑法不禁止有利于被告人的溯及既往。也即，事后的这个法律如果对被告人有利，也可以适用。

3.\CJKunderwave{严格}的罪刑法定

这是指禁止类推解释。类推解释因为既违反了民主主义，又违反了预测可能性原理，

% Source PDF physical page 21; see audit/source_coverage.json.
所以应当被禁止。例如，本节开头“串通拍卖案”中，将拍卖解释为投标，属于类推解释，应当被禁止，因此对狗蛋不能定串通投标罪。

[注意］刑法不禁止有利于被告人的类推解释。（2020年试题）

4.\CJKunderwave{确定}的罪刑法定

这是指罪刑规范应当明确、适当。

（1）明确性要求

[注意］简单罪状不违反明确性要求。简单罪状，是指仅规定了罪名，没有具体描述犯罪特征。例如，《刑法》第232条只规定“故意杀人的”，没有详细规定故意杀人罪的定义。之所以没有详细描述，是因为没有必要，一般人都知道什么是故意杀人。（2006年第1题)

（2）禁止绝对不定刑及绝对不定期刑

绝对不定刑，是指刑法条文只规定“犯……罪，判处刑罚”，不规定刑罚的种类；绝对不定期刑，是指刑法条文只规定“犯..罪，判处有期徒刑”，不规定具体刑期。因为绝对不定刑和绝对不定期刑都违反了预测可能性原理，应当被禁止。因此，没有犯罪就没有刑罚。反过来，没有明确的刑罚就没有犯罪。

（3）禁止处罚不当罚的行为

这是刑法谦抑性、补充性的要求。由于刑罚是最严厉的制裁措施，只有在其他法律已经无济于事、无法规制的情况下，才可以适用刑法，刑法的启动应当保持谦卑姿态。刑法是其他法律的补充法、保障法。

[注意］民法（及行政法）与刑法的关系。

二者不是对立排斥、非此即彼的关系，而是低位阶与高位阶的关系。\footnote{这个位阶关系，是从适用顺序而言，而非指效力等级。二者效力等级是相同的，都是全国人大制定的基本法律。}也即，不是A与-A的关系，而是A与A+B的关系。不能因为一个行为符合民法上的要件，就以此认定该行为不构成犯罪。

例如，甲将乙打成轻微伤，只是民法上的侵权行为；打成轻伤或重伤，既是民法上的侵权行为，又是刑法上的故意伤害行为，此时定故意伤害罪，然后附带民事诉讼。这是因为高位阶法优于低位阶法。

实务中常犯的错误：例如，对于“许霆案（利用ATM机出错非法占有银行资金）”有人认为许霆无罪，无罪的理由是许霆的行为属于民法上的不当得利。又如，某省高级法院曾颁布指导意见：肇事司机不逃逸，救助被害人，保护现场，属于履行行政法上的义务，不能认定为自首，若逃逸后又自动投案，则可认定为自首。这种错误意见的根源在于将行政法与刑法对立排斥起来。

（4）禁止不均衡、残虐的刑罚

禁止不均衡的刑罚，是指刑罚应与罪行相适应，重罪重判，轻罪轻判。这表明广义的罪刑法定原则可以包含罪刑相适应原则。禁止残虐的刑罚是人道主义的要求。

\subsection*{二、罪刑相适应原则}

第5条刑罚的轻重，应当与犯罪分子所犯罪行和承担的刑事责任相适应。

% Source PDF physical page 22; see audit/source_coverage.json.
\InsertDiagram{figures/manual_p022_1.tex}

例如，对自首之所以从宽处罚，是因为人身危险性降低了。对“入户抢劫”之所以适用更重的法定刑，是因为比普通抢劫罪多了一项法益侵害（非法侵入住宅）。

在考察不同犯罪人的人身危险性时，不能“一刀切”，应当具体地、个别化地判断，在追求一般正义的同时注意实现个别正义。这便是刑罚个别化的理念（2014年第1题）。

例如，甲乙都盗窃他人财物，甲是为了赌博，乙是为了给母亲治病。二人的人身危险性有所不同，所处刑罚应有所区别。

\OriginalPanel{典型真题}

关于社会主义法治理念与罪刑法定原则的关系有以下观点：

①罪刑法定的思想基础是民主主义与尊重人权主义，具备社会主义法治理念的本质属性

②罪刑法定既约束司法者，也约束立法者，符合依法治国理念的基本要求

③罪刑法定的核心是限制国家机关权力，保障国民自由，与执法为民的理念相一致

④罪刑法定是依法治国理念在刑法领域的具体表现

关于上述观点的正误，下列哪一选项是正确的？（2013年·卷二·2题）\footnote{[答案]D。}

A.第①句正确，第②③④句错误

B.第①③句正确，第②④句错误

C.第①②③句正确，第④句错误

D.第①②③④句均正确

\section{刑法的效力}

[问题]朝鲜公民金正男（领导人金正恩的哥哥）在马来西亚被越南籍女杀手杀死。哪些国家的刑法可以管辖本案？（金正男案）

\subsection*{一、刑法的空间效力}

\subsubsection*{（一）在中国境内犯罪：属地管辖原则}

第6条凡在中华人民共和国\CJKunderwave{领域内}犯罪的，除法律有特别规定的以外，都适用本法。

凡在中华人民共和国\CJKunderwave{船舶或者航空器}内犯罪的，也适用本法。

犯罪的\CJKunderwave{行为或者结果}有一项发生在中华人民共和国领域内的，就认为是在中华人民共和国领域内犯罪。

% Source PDF physical page 23; see audit/source_coverage.json.
1．“我国领域”包括领土、领水、领空\footnote{属地管辖的例外：法律有特别规定。这主要包括三种情形：第一，享有外交特权和豁免权的外国人在中国境内犯罪，通过外交途径解决；第二，我国港澳台地区适用本地区刑法；第三，民族自治区人大根据本民族特点，针对刑法部分条文制定了变通或补充规定，适用该规定。}

（1）旗国主义。悬挂我国国旗的航空器与船舶，不论停放何处，都属于我国领域。这是旗国主义的体现。（2005年第56题）

(2）例外。国际列车、国际长途汽车不属于我国领域。（2007年第51题）

2.属地管辖原则之“地”，既包括行为地，也包括结果地，二者只要具备其一即可。

（1）\CJKunderwave{犯罪行为}，从共同犯罪上分，包括实行行为、教唆行为和帮助行为；从犯罪形态上分，包括预备行为和实行行为。上述行为中，只要有一项行为发生在国内，其他相关行为即使发生在国外，也认为是在我国领域内犯罪。\footnote{注意：教唆、帮助行为在国内，实行行为在国外。（1）情形一：根据当地国刑法，实行行为无罪（我国刑法规定为犯罪），则教唆行为、帮助行为也无罪，我国刑法不需追究。例如，甲在国内教唆美国人乙在拉斯维加斯开设赌场。由于乙的行为在美国无罪（我国刑法有开设赌场罪），那么甲的教唆行为也无罪，我国刑法不需追究。（2）情形二：根据当地国刑法及我国刑法，实行行为有罪（只是我国刑法无法追究），则教唆、帮助行为有罪，我国刑法需要追究。例如，甲在国内教唆美国人乙在日本杀害日本人丙。日本及我国刑法都规定有杀人罪，乙的行为有罪（只是我国刑法无法追究，也不需要追究），因此甲的教唆行为也有罪，我国刑法予以追究。}（2007年第51题）

（2）\CJKunderwave{犯罪结果}，从共同犯罪上分，包括整体结果和部分共犯人的结果；从犯罪形态上分，包括实害结果和危险结果。例如，甲从美国给我国境内的乙邮寄炸弹，乙收到炸弹后，炸弹爆炸了（实害结果），或者炸弹被海关查获没爆炸（危险结果）。上述结果中，只要有一个结果发生在国内，就认为是在我国领域内犯罪。\footnote{注意1：犯罪地不包括单纯过境地。例1，甲从日本寄毒药给身在中国的乙，乙服毒后来到法国，在法国死亡。由于在我国国内发生了紧迫危险，我国作为危险结果地能够管辖此案。例2，甲从日本向法国的乙寄送毒药，运送该毒药的飞机经过我国领空，在我国机场停留了两个小时。由于在我国国内没有发生紧迫危险，我国不属于危险结果地，不管辖此案。

注意2：跨境网络犯罪的犯罪地，包括：网站服务器所在地，网络接人地，网站建立者、管理者所在地，犯罪活动参与者所在地，被害人所在地。例1，缅甸人甲在缅甸建立赌博网站，中国境内的人在该网站参赌或者做代理，那么中国境内属于该开设赌场罪的犯罪地。例2，甲在马来西亚向中国境内的人打电话诈骗，或者实施网络诈骗，那么中国境内属于该诈骗罪的犯罪地。}

\subsubsection*{（二）在中国境外犯罪}

1.属人管辖原则（第7条）：我国公民在境外犯罪

第7条中华人民共和国公民在中华人民共和国\CJKunderwave{领域外}犯本法规定之罪的，适用本法，但是按本法规定的\CJKunderwave{最高刑为三年以下有期徒刑}的，\CJKunderwave{可以不予追究}。

中华人民共和国\CJKunderwave{国家工作人员和军人}在中华人民共和国领域外犯本法规定之罪的，适用本法。

适用我国刑法的条件：第一，我国公民在国外犯我国刑法规定的犯罪，原则上适用我国刑法，犯轻罪的（最高刑在3年或以下），可以不予追究，这意味着也可以追究；第二，若为国家工作人员和军人犯罪的，一律追究。

2.保护管辖原则（第8条）：外国人在境外犯罪

第8条外国人在中华人民共和国领域外对中华人民共和国国家或者公民犯罪，而按本法规定的\CJKunderwave{最低刑为三年以上有期徒刑}的，可以适用本法，但是\CJKunderwave{按照犯罪地的法律不受处}\CJKunderwave{罚的除外}。

% Source PDF physical page 24; see audit/source_coverage.json.
适用我国刑法的条件：第一，针对我国国家或公民的犯罪；第二，行为触犯的是重罪（最低刑在3年或以上）；第三，双重犯罪原则（犯罪地的法律也认为是犯罪）。（2007年第51题)

例如，本节开头“金正男案”中，朝鲜依据保护管辖，马来西亚依据属地管辖，越南依据属人管辖。三个国家都可以管辖，此时属地管辖优先，因为主要证据在马来西亚。

3.普遍管辖原则（第9条）：国际犯罪

第9条对于中华人民共和国缔结或者参加的国际条约所规定的罪行，中华人民共和国在所承担条约义务的范围内行使刑事管辖权的，适用本法。

适用我国刑法的条件：第一，必须是危害人类共同利益的犯罪，主要有：劫持航空器、跨国贩毒、跨国拐卖人口、海盗、种族灭绝、洗钱、恐怖活动等；第二，我国缔结或参加了公约；第三，我国刑法将这种行为也规定为犯罪；第四，犯罪人出现在我国领域内。注意：具体适用法律时，适用我国刑法，而非已缔结或参加的国际公约。（2004年第56题)

\InsertDiagram{figures/manual_p024_1.tex}

\subsection*{二、刑法的时间效力}

\InsertDiagram{figures/manual_p024_2.tex}

我国刑法关于时间效力的原则是：\CJKunderwave{从旧兼从轻}。这是指，原则上适用旧法（行为时的法律），但适用新法有利于被告人时，适用新法。也即，原则上禁止法律溯及既往，例外是不禁止有利于被告人的溯及既往。\footnote{现行刑法于1997年10月1日生效。《刑法修正案（八）》于2011年5月1日生效。《刑法修正案（九）》于2015年11月1日生效。《刑法修正案（十）》于2017年11月4日生效。《刑法修正案（十一）》于2021年3月1日生效。《刑法修正案（十二）》于2024年3月1日生效。}

例如，甲是某民企的经理，在2023年利用职务便利，为亲友非法牟利。2024年4月案发。《刑法修正案（十二）》生效前，为亲友非法牟利罪的犯罪主体仅限于国有单位的人。2024年3月1日，《刑法修正案（十二）》生效，该罪的犯罪主体包括了民企的人。

% Source PDF physical page 25; see audit/source_coverage.json.
按照行为时的规定，甲不构成为亲友非法牟利罪。按照新规定，甲构成该罪。对此，应当从旧，也即按照行为时的规定。

1.适用的案件

“从旧兼从轻”原则适用的对象是未决犯（未判决的案件），也即，案件发生在旧法时，审判发生在新法时。对于已决犯（已判决的案件）则不存在是否溯及既往的问题。

\CJKunderwave{注意}：按照审判监督程序（再审程序）重新审判的案件，适用行为时的法律。例如，甲在1990年对妇女要流氓，按照当时刑法，符合流氓罪构成要件，被判流氓罪，1995年刑满释放。1997年新刑法（现行刑法）生效，删除流氓罪。甲申诉，要求再审。对该案再审时，应适用行为时的法律。只要行为时的事实认定和法律适用没有问题，就应维持原判。

2.司法解释的时间效力\footnote{参见2001年12月17日最高人民法院、最高人民检察院《关于适用刑事司法解释时间效力问题的规定》。}（2017年第1题）

\InsertDiagram{figures/manual_table_p025_1.tex}
\SourceFootnote{\textbf{表注②}\quad 司法解释具有溯及力，似乎违反了罪刑法定原则，实际上不违反。这是因为，司法解释不是立法条文，只是对立法条文的含义的解释，因此按照司法解释处理，不违反罪刑法定原则。}

